\chapter{Anatomy of \texorpdfstring{\texttt{(aux kanren micro)}}{(aux kanren micro)}}\label{ch:micro}

% ------------------------------------------------------------------------------------------------
\section*{What and why}

The chapters so far studied what a higher-order logic programming language \emph{is}: the uniform
proofs of hereditary Harrop formulas (\cref{ch:hh}), λ-terms as data (\cref{ch:lambda,ch:lambdatree}),
higher-order pattern unification (\cref{ch:hounif}), and, in \cref{ch:lkanren}, how λKanren grafts
these features onto a miniKanren~\cite{MaFriedman2021,MaLuFriedman2020}. This chapter opens the
\emph{proposal} part of the document. Before deciding how \micro{} should grow, we need an exact
picture of the engine as it is today: what a variable is, what a state contains, how unification
and the constraint stores interact, how the search interleaves, and how answers are printed. The
point is not to explain the code to its author but to put next to each other the facts that matter
when the engine is extended: which invariants a new feature must preserve, which parts can be
reused from outside, and which parts are closed.

The engine is a μKanren in the sense of Hemann and Friedman~\cite{HemannFriedman2013} (a goal maps a
state to a stream of states) with a cKanren-like constraint store holding disequalities, disjoint
type tags and \emph{absento} constraints, in the style of \emph{The Reasoned
Schemer}~\cite{FriedmanEtAl2018}. All references \texttt{:N} are line numbers of
\texttt{src/aux.kanren.micro.scm} (713 lines) at commit \texttt{f53bdcb}. Every Scheme snippet in this
chapter was run on that engine; its output is pasted verbatim below it. The snippets import
\scm{(aux base)}, \scm{(aux stream)}, \scm{(aux fds sbral)} and \micro.

\begin{table}[h]
\centering\small
\caption{A map of \texttt{aux.kanren.micro.scm}.}\label{tab:micro-map}
\begin{tabular}{@{}lll@{}}
\toprule
lines & content & section \\
\midrule
\texttt{3–26}    & module header, explicit export list, imports & \\
\texttt{28–59}   & sentinel, tags, variables & \cref{sec:micro-terms} \\
\texttt{60–107}  & the state record, sbral helpers & \cref{sec:micro-state,sec:micro-subst} \\
\texttt{110–204} & occurs check and binding, walks, unification & \cref{sec:micro-subst,sec:micro-unify} \\
\texttt{206–326} & reification and projection & \cref{sec:micro-reify} \\
\texttt{330–523} & D/A/T verification and subsumption & \cref{sec:micro-stores} \\
\texttt{525–629} & goals, combinators, constraint goals & \cref{sec:micro-search} \\
\texttt{631–684} & aggregation & \cref{sec:micro-aggr} \\
\texttt{686–713} & \texttt{define-relation}, runners & \cref{sec:micro-search,sec:micro-reify} \\
\bottomrule
\end{tabular}
\end{table}

% ------------------------------------------------------------------------------------------------
\section{Terms and variables}\label{sec:micro-terms}

There is no term datatype: a term is any Scheme value, and the engine gives structure to some of
them.

\begin{definition}[Terms of \micro]\label{def:micro-term}
A \emph{term} is one of: a logic variable; an atom (symbol, boolean, \scm{'()}, character, number,
string, or any other Scheme object); a pair \scm{(t . u)} of terms; a vector of terms; an instance of a
CHICKEN record type (\scm{define-record}) whose slots are terms. Vectors and record instances
together are the \emph{structures} (\scm{μkanren-structure?}, \texttt{:110}).
\end{definition}

A logic variable is a record with a single integer slot (\texttt{:41}):
\begin{scheme}
(define-record μkanren-var index)
(define (μkanren-var-working? α) (and (μkanren-var? α) (<= 0 (μkanren-var-index α))))
\end{scheme}
The sign of the index separates two kinds of variable:
\begin{itemize}
\item a \emph{working} variable has index $i \ge 0$; it lives during the search, its index is its
  slot in the substitution, and it prints as \scm{_i};
\item a \emph{reified} variable has index $-1,-2,\dots$; it exists only in answers, is created only by
  reification (\cref{sec:micro-reify}), and prints as a Greek letter \scm{α}, \scm{β}, \dots,
  \scm{ω}, then \scm{α1}, \dots (\texttt{:52–59}).
\end{itemize}
The helper \scm{μkanren-var-index>=0} (\texttt{:50}) maps both kinds onto $\mathbb{N}$:
$i\mapsto i$ for working variables, $-1\mapsto 0$, $-2\mapsto 1$, \dots{} for reified ones.

\begin{scheme}
(write (map make-μkanren-var '(0 1 7 -1 -2 -24 -25))) (newline)
\end{scheme}
\begin{answers}
(_0 _1 _7 α β ω α1)
\end{answers}

Indices are allocated by \scm{freshª} from a counter in the state (\cref{sec:micro-search}), so along
one branch of the search they increase monotonically: \emph{the index of a variable is its age},
lower meaning older. This fact, invisible to a user, is the most useful one for the λProlog
extension.

\begin{pitfall}[A variable is a record]
A \scm{μkanren-var} is itself a \scm{record-instance?}. Every structural walker of the engine
therefore tests \scm{μkanren-var?} \emph{before} the generic record case (e.g.\ \texttt{:117} before
\texttt{:124}, and the explicit \scm{(not (μkanren-var? w*))} at \texttt{:190}). Any new walker,
such as a scope check, must do the same, or it will descend into a variable as if it were a
two-slot structure.
\end{pitfall}

\begin{kanren}[Indices as levels]
λProlog's GENERIC rule reduces $\seq{\Sig}{\Prog}{\forall x\,G}$ to
$\seq{c{:}\tau,\Sig}{\Prog}{G[c/x]}$ for a token $c$ not in $\Sig$ \book{39}, and the query
\lp{sigma x\ pi y\ x = y} must fail \book{85--86}: with logic variables, a variable created
\emph{before} $c$ must never be bound to a term containing it. Implementations enforce this
with universe indices or levels attached to variables and constants~\cite{NadathurJayaramanKwon1995}.
In \micro{} the monotone index already is a level boundary: if the eigenvariable is created when the
counter is $n$, then ``created before $c$'' is exactly ``index $< n$''. The prototype of
\cref{sec:micro-forall} uses precisely this. Index order alone is not enough, though: a variable created
after the $\forall$ has exited must not see $c$ either, and lowering changes a variable's level after its
creation. The proposal therefore keeps a level per variable, in an \texttt{sbral} indexed like S and padded
only when the depth changes or the unifier creates a variable of another level (\cref{def:design-levels},
\cref{def:module-state}, \cref{sec:module-state}).
\end{kanren}

% ------------------------------------------------------------------------------------------------
\section{The state}\label{sec:micro-state}

\begin{definition}[State]\label{def:micro-state}
A state is a record of six fields (\texttt{:69}):
\begin{scheme}
(define-record μkanren-state vars-count S D A T tags)
(define μkanren-state-empty (make-μkanren-state 0 empty/sbral empty/sbral empty/sbral empty/sbral '()))
\end{scheme}
\end{definition}

\begin{longtable}{@{}>{\raggedright\arraybackslash}p{0.14\fullwidthlen}>{\raggedright\arraybackslash}p{0.36\fullwidthlen}>{\raggedright\arraybackslash}p{0.42\fullwidthlen}@{}}
\toprule
field & content & written by / read by \\
\midrule\endhead
\texttt{vars-count} & the next fresh index, i.e.\ the number of variables introduced on this branch
  & written and read by \scm{freshª} (\texttt{:532–537}) \\
\texttt{S} & the substitution: a skew-binary random-access list (sbral) of terms or of the sentinel
  \scm{*μkanren-unbound*} & written only by \scm{μkanren-state-update} (\texttt{:113–134}); read by every walk \\
\texttt{D} & disequality store: a list of \emph{prefixes}, each a list of \scm{(var . term)} pairs
  & \scm{≠°} (\texttt{:554–567}), re-verification after \scm{=°} \\
\texttt{A} & type store: a list of \scm{(var . tag)}; the tags are pairwise disjoint
  & \scm{symbol°}, \scm{number°}, user tags; verified after \scm{=°} \\
\texttt{T} & absento store: a list of \scm{(term . tag)}, the term a variable or a compound
  & \scm{absent°} (\texttt{:617–624}); verified after \scm{=°} \\
\texttt{tags} & tags asserted through A & written at \texttt{:481}, read by nobody in the engine (a dead field) \\
\bottomrule
\end{longtable}

Only S is a real sbral; D, A and T are plain lists (the \scm{empty/sbral} in the empty state is just
\scm{'()}). States are immutable and copied by functional update with keyword arguments
(\texttt{:76–83}),
\begin{scheme}
(μkanren-state-with s D: (append D* D))          ; a copy of s with a new D
\end{scheme}
except at the two hottest sites, \scm{μkanren-state-update} and \scm{freshª}, which call the
positional constructor directly. A positional destructuring macro \scm{μkanren-state-match} of fixed
arity six (\texttt{:96–103}) is used by six goals (\texttt{:130}, \texttt{:515}, \texttt{:533},
\texttt{:540}, \texttt{:555}, \texttt{:618}).

To see states, we use a small printer (not part of the engine) and run goals on the empty state:
\begin{scheme}
(define (dump s)   ; the fields of a state on one line; a tag is shown by its name
  (for-each (λ (k v) (display k) (display " ") (write v) (display "  "))
            '(vars-count S D A T)
            (list (μkanren-state-vars-count s)
                  (sbral->list (μkanren-state-S s))
                  (μkanren-state-D s)
                  (map (λ (p) (cons (car p) (μkanren-tag-name (cdr p)))) (μkanren-state-A s))
                  (map (λ (p) (cons (car p) (μkanren-tag-name (cdr p)))) (μkanren-state-T s))))
  (newline))
(define (states g) (§->list (g μkanren-state-empty)))   ; every answer state of g

(for-each dump (states (fresh° (q x y) (=° q (list x y)) (≠° x 1) (symbol° y) (absent° 'k x))))
\end{scheme}
\begin{answers}
vars-count 3  S ((_1 _2))  D (((_1 . 1)))  A ((_2 . μkanren-tag/sym87))  T ((_1 . k))
\end{answers}
The three variables \scm{q}, \scm{x}, \scm{y} are \scm{_0}, \scm{_1}, \scm{_2}; S holds one binding
(\scm{_0} is bound to the list \scm{(_1 _2)}); each store holds one constraint. The numeric suffix
of a tag name such as \scm{μkanren-tag/sym87} is a \scm{gensym} counter (\texttt{:33–34}): it identifies
the tag but its value depends on the program that loaded the module, so the same tag prints as
\scm{μkanren-tag/sym98} in another script of this chapter.

\begin{proposition}[The eq?-invariant]\label{prop:micro-eq}
States are compared by identity: \scm{(define (μkanren-state-equal? s s*) (eq? s s*))} (\texttt{:73}).
Unification returns \emph{the very same state object} when it adds no binding, and a freshly
allocated one otherwise (by inspection of \texttt{:160–181}: the non-binding cases call
\scm{return}, which closes over the input state, and only \scm{μkanren-state-update} allocates).
\end{proposition}

This is not an optimisation but semantics: every entailment test of the engine reads ``unification
gave back \scm{s}'' as ``the two terms are already equal''. The users are \scm{≠°} (\texttt{:559}: the
terms are already equal, so the disequality fails), \scm{=°} (\texttt{:544}: nothing new, so skip
re-verification), the re-verification of one disequality (\texttt{:460}), the absento test
\scm{μkanren-tag-unify} (\texttt{:394}) and subsumption (\texttt{:304}).

\begin{pitfall}[Preserving the eq?-invariant]
Any extension that runs inside unification, or after it, must return the input state itself when it
changes nothing. A hook that rebuilds the state with \scm{μkanren-state-with}, which always allocates,
on a no-op unification would make \scm{(≠° 1 1)} look satisfiable.
\end{pitfall}

\paragraph{How the record has grown.} The state record was widened in place several times: from
\scm{(substitution counter)} to \scm{(S D A T)} (commit \texttt{61bb868}), then \texttt{vars-count}
(\texttt{1e7994e}), then \texttt{tags} (\texttt{d118ddd}). Each step edited micro and the positional
state constructions of \texttt{test/microkanren.scm}. No sibling module (arith, date, show, unionfind)
ever builds a state: every store so far was added \emph{inside} micro. This is the precedent the
proposal of \cref{ch:design} has to reckon with.

% ------------------------------------------------------------------------------------------------
\section{The substitution and the walk}\label{sec:micro-subst}

\paragraph{S as an sbral.} \texttt{(aux fds sbral)} implements Okasaki's skew-binary random-access
lists: \scm{cons/sbral}, \scm{car/sbral}, \scm{cdr/sbral} in $O(1)$, \scm{sbral-ref} and
\scm{update/sbral} in $O(\log n)$, \scm{length/sbral} in $O(\log n)$. S stores the \emph{most recent
variable first}: the slot of the working variable of index $i$ sits at position $\ell-1-i$, where
$\ell$ is the length of S (\scm{μkanren-var-index/sbral}, \texttt{:86}). Extending S at the front
with \scm{cons/sbral} therefore keeps the positions of older variables valid when counted from the
end.

\paragraph{Lazy extension.} \scm{freshª} does not touch S. Binding the variable of index $i$
first pads S with the sentinel \scm{*μkanren-unbound*} until it has a slot $i$
(\scm{μkanren-var-extend/sbral}, \texttt{:89–94}). Hence
$\ell = 1+\max\{\text{bound index}\}$, and a variable whose index is $\ge \ell$ is unbound
\emph{without} having a slot.

\begin{scheme}
(for-each dump (states (fresh° (a b c d) (=° c 5))))
\end{scheme}
\begin{answers}
vars-count 4  S (5 #<aux.kanren.micro#μkanren-unbound> #<aux.kanren.micro#μkanren-unbound>)  D ()  A ()  T ()
\end{answers}
Four variables were introduced but S has three slots: \scm{c} (index 2) is bound to 5 at the front,
\scm{a} and \scm{b} hold the sentinel, \scm{d} (index 3) lies beyond S.

\begin{definition}[Walk]\label{def:micro-walk}
Let $\ell$ be the length of S. The \emph{walk} of a term $t$ (\scm{μkanren-state-find/SL},
\texttt{:139–145}) is computed by the loop $F(\beta_0,\beta)$, started with $\beta_0=\beta=t$:
\begin{enumerate}
\item if $\beta$ is a variable with $\mathit{idx}(\beta) \ge \ell$ (with $\mathit{idx}$ $=$ \scm{μkanren-var-index>=0}), answer $\beta$ (unbound, no slot);
\item if $\beta$ is a working variable, continue with $F(\beta, S[\ell-1-i])$;
\item if $\beta$ is the sentinel, answer $\beta_0$, the last variable visited;
\item otherwise answer $\beta$.
\end{enumerate}
\end{definition}

Case~4 catches non-variables but also \emph{reified} variables, which are not working: to the walk
a reified variable is a constant. Bindings are triangular (S maps a variable to a term that may
contain bound variables) and the walk never compresses paths. The variant with explicit \scm{S} and
\scm{l} hoists the $O(\log n)$ \scm{length/sbral} out of deep traversals; \scm{μkanren-state-find}
(\texttt{:147–148}) computes it once.

\paragraph{Deep walks.} Three exported procedures resolve a term completely:
\begin{itemize}
\item \scm{μkanren-state-find/value} (\texttt{:183–191}) returns a Scheme value, rebuilding pairs with
  \scm{cons}, vectors with \scm{list->vector} and records with \scm{(apply make-record-instance slots)};
  unbound variables stay \scm{μkanren-var} records;
\item \scm{μkanren-state-find/values} (\texttt{:632}) maps it over a list (the key of an aggregation group);
\item \scm{μkanren-state-find/repr} (\texttt{:193–204}) returns an \emph{expression} that rebuilds the
  value: a variable becomes its symbol, a symbol \scm{'sym}, a pair \scm{(cons a d)}, a vector
  \scm{(vector ...)}, a record \scm{(make-record-instance 'type ...)}.
\end{itemize}

\begin{pitfall}[Deep walks rebuild records]
Because \scm{find/value} rebuilds every record, the value it returns is \scm{equal?} but not
\scm{eq?} to the record that was unified. Here a record \scm{c} of the type \scm{eigen} of
\cref{sec:micro-forall} is bound to \scm{q} and read back:
\begin{scheme}
(let* ((c (make-eigen 99)) (r (car (§->list ((fresh° (q) (=° q c)) μkanren-state-empty))))
       (v (μkanren-state-find/value (make-μkanren-var 0) r)))
  (pretty-print (list (eq? v c) (equal? v c) (eigen? v))))
\end{scheme}
\begin{answers}
(#f #t #t)
\end{answers}
Constructed terms must be compared with predicates and \scm{equal?} (or by an id slot), never with
\scm{eq?}; \cref{sec:micro-forall} shows a scope check that goes wrong exactly this way.
\end{pitfall}

% ------------------------------------------------------------------------------------------------
\section{Unification}\label{sec:micro-unify}

\subsection{Binding and the occurs check}

Exactly one procedure creates a binding (\texttt{:113–134}, abridged):
\begin{scheme}
(define (μkanren-state-update α v s assocs #!key (unsafe #f))
  (define (occur? v*) ...)                  ; does α occur in the walk of v*?
  (and
    (or unsafe (not (occur? v)))
    (μkanren-state-match ((vc S D A T tags) s)
      (let* ((S* (μkanren-update/sbral α v S))
             (s* (make-μkanren-state vc S* D A T tags))
             (assocs* (cons `(,α . ,v) assocs)))
        `(,s* ,assocs*)))))
\end{scheme}
It answers \scm{#f} if $\alpha$ occurs in $v$ (the occurs check walks $v$, its pairs and its structures
through S), otherwise the pair \scm{(s* assocs*)}. The list \scm{assocs} accumulates \emph{the bindings
made by the current unification} as \scm{(var . term)} pairs: this is cKanren's ``prefix''. The
\scm{unsafe:} key is never passed, so the occurs check is always on. The procedure is not exported;
it is the only place where a binding can be refused before it enters the state, e.g.\ for violating a
scope; a check run right after each \scm{=°}, before the next goal, works as well (the route the proposal
takes, \cref{sec:design-hooks}).

\subsection{The unification algorithm}

\begin{definition}[Unification, \texttt{:160–181}]\label{def:micro-unify}
\scm{(μkanren-state-unify u v s assocs)} walks both sides to $u^*$, $v^*$ and tries these cases
\emph{in order}:
\begin{enumerate}
\item $u^* \mathrel{\texttt{eq?}} v^*$: success, nothing bound;
\item two strings that are \scm{string=?}: success;
\item two numbers that are \scm{μkanren-number-eqv?} (\texttt{:153–158}: \scm{eqv?}, but two NaN
  flonums are equal and non-real numbers are compared part by part): success;
\item two working variables of indices $i, j$: if $i=j$ success; if $i<j$ bind $u^* := v^*$, else
  $v^* := u^*$, i.e.\ \emph{the older variable is bound to the newer one};
\item one working variable: bind it to the other side;
\item two structures: unify the lists of their slots (\scm{μkanren-structure->list}, \texttt{:111}),
  where \scm{record->vector} puts the record's type name first;
\item two pairs: unify the cars, then the cdrs in the resulting state;
\item otherwise fail with \scm{#f}.
\end{enumerate}
The result is \scm{(s* assocs*)} or \scm{#f}. With \scm{assocs} $=$ \scm{'()}, the second component is
exactly the list of new bindings.
\end{definition}

Read as rules, the algorithm is the standard first-order syntactic one with occurs check; what is
specific to \micro{} is (a) the orientation of variable–variable bindings, (b) structures, (c) the
treatment of atoms. The exported procedure can be called on a raw state:
\begin{scheme}
(define q (make-μkanren-var 0))
; a state where _0 exists
(define s0 (μkanren-state-with μkanren-state-empty vars-count: 1))

(define r (μkanren-state-unify q '(1 2) s0 '()))
(pretty-print (cadr r))                                         ; the new bindings
(define r2 (μkanren-state-unify q '(1 2) (car r) '()))
; nothing new: the very same state
(pretty-print (list (eq? (car r2) (car r)) (cadr r2)))
(pretty-print (μkanren-state-unify q (list 1 q) s0 '()))        ; occurs check
\end{scheme}
\begin{answers}
((_0 1 2))
(#t ())
#f
\end{answers}
The first line is the one-element prefix \scm{((_0 . (1 2)))}, printed as Scheme prints it; the second
shows \cref{prop:micro-eq} at work; the third is the occurs check.

\paragraph{Orientation.} Unifying two fresh variables binds the older to the newer:
\begin{scheme}
(for-each dump (states (fresh° (a b) (=° a b))))
\end{scheme}
\begin{answers}
vars-count 2  S (_1)  D ()  A ()  T ()
\end{answers}
S has one slot, the one of \scm{_0}, and it holds \scm{_1}. A level discipline in the style of
λProlog would prefer the opposite orientation for a variable–variable pair (bind the younger variable
to the older one, so that such a binding never moves a variable out of its scope). Flipping the
orientation would not be enough, though: in \lp{sigma X\ pi y\ sigma Z\ (X = f Z, Z = y)} the binding
$X := f\,Z$ makes the young $Z$ reachable from the old $X$, so a real level discipline must also lower
the level of every variable of a term bound to an older variable (this is the shape of the fourth
query of \cref{sec:micro-forall}, with \scm{list} for $f$). With the orientation of \micro{} an
old variable can moreover be aliased to a young one, and a scope check must look through such chains
(\cref{sec:micro-forall}).

\paragraph{Structures and atoms.} Records unify as the list of their slots, so any
\scm{define-record} is a new term constructor for free; but a record and a vector with the same slot
list unify too. Numbers are compared with \scm{eqv?}, so an exact and an inexact number never unify.
\begin{scheme}
(define-record dbi n)
(pretty-print (record->vector (make-dbi 3)))
(pretty-print (°->list #f (fresh° (q) (=° (make-dbi q) (vector 'dbi 7)))))
(pretty-print (°->list #f (fresh° (q) (=° (vector 1 q) '(1 2)))))
(pretty-print (°->list #f (fresh° (q) (=° (/ 0. 0.) (/ 0. 0.)) (=° "ab" (string #\a #\b)) (=° q 'ok))))
(pretty-print (°->list #f (fresh° (q) (=° 1 1.0) (=° q 'ok))))
\end{scheme}
\begin{answers}
#(dbi 3)
((λ () 7))
()
((λ () 'ok))
()
\end{answers}
In order: the slot list of a record starts with its type name; a \scm{dbi} record unifies with the
vector \scm{#(dbi 7)}; a vector never unifies with a list; NaN equals NaN and equal strings unify
although they are distinct objects; \scm{1} and \scm{1.0} do not unify. Anything else (procedures,
bytevectors, \dots) unifies only with itself under \scm{eq?}.

\begin{pitfall}[Procedures are opaque]
A Scheme closure is an atom compared with \scm{eq?}. Higher-order abstract syntax with real Scheme
procedures (the λProlog \lp{abs x\ ...} rendered as \scm{(λ (x) ...)}) cannot be unified, walked,
constrained or reified. λ-terms must be first-order data, e.g.\ de Bruijn records such as \scm{dbi}
above (\cref{ch:lambdatree}).
\end{pitfall}

\begin{kanren}[Where higher-order unification would go]
The \scm{cond} of \cref{def:micro-unify} is closed: there is no dispatch table and no hook. A new term
kind with \emph{syntactic} equality needs nothing (it is a record); a term kind with a different
equality, such as $\Llam$ flexible terms $F\,x_1\dots x_n$ solved by pattern unification
(\cref{ch:hounif}) or nominal terms with suspended swappings~\cite{ByrdFriedman2007,UrbanPittsGabbay2004},
needs either a new clause here (before case~6) or a goal that solves the problem outside and emits
ordinary bindings (\cref{sec:micro-noedit}). The choice between them is made in \cref{ch:design}.
\end{kanren}

% ------------------------------------------------------------------------------------------------
\section{The constraint stores and their re-verification}\label{sec:micro-stores}

Unification only binds; the three stores keep \emph{negative} and \emph{type} information, and must
be re-checked whenever a binding is made.

\subsection{D: disequalities}

\begin{definition}[Prefix]\label{def:micro-prefix}
A \emph{prefix} is a list $d = ((x_1 . t_1) \dots (x_k . t_k))$ of variable–term pairs; it denotes
the constraint $\neg(x_1 = t_1 \wedge \dots \wedge x_k = t_k)$. The store D is a list of prefixes,
read as their conjunction.
\end{definition}

The goal \scm{(≠° u v)} (\texttt{:554–567}) runs a \emph{trial} unification of $u$ and $v$ with an
empty \scm{assocs} and reads the result:
\begin{itemize}
\item \scm{#f}: $u$ and $v$ can never be equal, so the constraint is dropped and the goal succeeds
  with the state unchanged;
\item the very same state (\cref{prop:micro-eq}): $u$ and $v$ are already equal, so the goal fails;
\item a new state with bindings $d$: the goal succeeds with the \emph{old} state extended by the
  prefix $d$ in D, unless a type tag of A or an absento of T makes $d$ vacuous (subsumption,
  below).
\end{itemize}
A tuple disequality gives one prefix with several pairs. After a later \scm{=°}, a prefix is
re-verified by trial unification of its pairs (\scm{μkanren-verify-D+}, \texttt{:458–461}): if the
trial fails the prefix is satisfied forever and dropped; if it adds nothing the prefix is violated and
the \scm{=°} fails; otherwise the prefix is replaced by its normalised version, the new trial bindings.

\begin{tutorial}[The life of a tuple disequality]
We assert $(x,y) \ne (1,2)$ and then bind $x$, watching D.
\begin{scheme}
(for-each dump (states (fresh° (q x y) (=° q (list x y)) (≠° (list x y) '(1 2)))))
(for-each dump (states (fresh° (q x y) (=° q (list x y)) (≠° (list x y) '(1 2)) (=° x 1))))
(for-each dump (states (fresh° (q x y) (=° q (list x y)) (≠° (list x y) '(1 2)) (=° x 3))))
(write (states (fresh° (q x y) (=° q (list x y)) (≠° (list x y) '(1 2)) (=° x 1) (=° y 2)))) (newline)
\end{scheme}
\begin{answers}
vars-count 3  S ((_1 _2))  D (((_2 . 2) (_1 . 1)))  A ()  T ()
vars-count 3  S (1 (_1 _2))  D (((_2 . 2)))  A ()  T ()
vars-count 3  S (3 (_1 _2))  D ()  A ()  T ()
()
\end{answers}
\begin{enumerate}
\item The trial unification of \scm{(_1 _2)} with \scm{(1 2)} would bind both variables, so D holds the
  single prefix \scm{((_2 . 2) (_1 . 1))}: ``not both''.
\item After \scm{(=° x 1)}, \scm{_1} is touched; the trial now binds only \scm{_2}, and the prefix is
  normalised to \scm{((_2 . 2))}: the constraint has become $y \ne 2$.
\item After \scm{(=° x 3)} the trial fails: the prefix can never be violated and disappears.
\item After \scm{(=° x 1)} and \scm{(=° y 2)} the trial binds nothing: the prefix is violated and the
  goal has no answer state.
\end{enumerate}
\end{tutorial}

\paragraph{Incremental re-verification.} After an \scm{=°}, only the prefixes that mention a
variable just bound (on either side of a pair) are re-verified (\scm{μkanren-verify-D/touched},
\texttt{:472–478}); the \scm{assocs} returned by unification is exactly the list needed to decide
which prefixes are touched. An untouched prefix cannot have become violated: its keys are unbound
variables, and since none of them was just bound, its trial unification would still bind its first key
(the reasoning of the comment at \texttt{:466–471}). It may, however, have become satisfied through the
occurs check deep in its terms; projection runs one last full \scm{μkanren-verify-D} to drop
those (\texttt{:466–471}, \texttt{:317}).

\subsection{A: disjoint type tags}

A \emph{tag} is a record \scm{(name def pred)} (\texttt{:31}): an identity, the quoted source of a
predicate (printed in answers), and the predicate itself. Type tags are made with
\scm{define-μkanren-tag} and turned into goals by \scm{μkanren-make-tag-A} (\texttt{:498–504}); the
built-in ones are (\texttt{:611–615})
\begin{scheme}
(define-μkanren-tag (μkanren-tag/sym v) (symbol? v))
(define-μkanren-tag (μkanren-tag/num v) (number? v))
(define symbol° (μkanren-make-tag-A μkanren-tag/sym))
(define number° (μkanren-make-tag-A μkanren-tag/num))
\end{scheme}
On a bound term the goal just tests the predicate; on an unbound variable it extends A. Tags are
\emph{mutually exclusive}: \scm{μkanren-ext-A} (\texttt{:431–436}) answers \scm{#f} when the variable
already carries a different tag. Adding a tag also drops the prefixes of D that it makes vacuous
(\scm{μkanren-subsume}, \texttt{:370–371}: a pair \scm{(α . u)} whose $\alpha$ is tagged by a predicate
rejecting $u$).
\begin{scheme}
(for-each dump (states (fresh° (q) (≠° q 'a) (number° q))))
(write (states (fresh° (q) (number° q) (symbol° q)))) (newline)
\end{scheme}
\begin{answers}
vars-count 1  S ()  D ()  A ((_0 . μkanren-tag/num99))  T ()
()
\end{answers}
A number can never be the symbol \scm{a}, so the disequality is gone; and no term is both a number and a
symbol.

\subsection{T: absento}

\leavevmode\scm{(absent° datum term)} (\texttt{:617–624}; note the argument order) states that \emph{datum}, which
may be any Scheme value, atomic or compound, does not occur in \emph{term}. The store holds
\scm{(term . tag)} pairs where the tag is an ad-hoc tag whose name is the datum itself. The term is a
variable, or a compound that could still become the datum (it is then ``watched whole'' and
decomposed only when it no longer can, \scm{μkanren-verify-T+}, \texttt{:410–419}). When a watched
variable is bound, the constraint moves to the parts:
\begin{scheme}
(for-each dump (states (fresh° (q x y) (absent° 'k q) (=° q (list x y)))))
(write (states (fresh° (q x y) (absent° 'k q) (=° q (list x y)) (=° y 'k)))) (newline)
\end{scheme}
\begin{answers}
vars-count 3  S ((_1 _2))  D ()  A ()  T ((_2 . k) (_1 . k))
()
\end{answers}
When a variable carrying an absento gets a type tag, the absento is converted into a disequality
(e.g.\ \scm{symbol° α} plus \scm{absent° 'closure α} becomes the prefix \scm{((α . closure))}), by
\scm{μkanren-subsume-A} and \scm{μkanren-update-D/T} (\texttt{:340–356}, \texttt{:480–485}). The
conversion happens for \emph{any} tag, and because the subsumption of D by the new tag runs before it
(\texttt{:493}), a converted prefix that the tag makes vacuous stays in the state and is dropped only at
projection:
\begin{scheme}
(for-each dump (states (fresh° (q) (absent° 'closure q) (symbol° q))))
(for-each dump (states (fresh° (q) (absent° 'closure q) (number° q))))
(pretty-print (°->list #f (fresh° (q) (absent° 'closure q) (number° q))))
\end{scheme}
\begin{answers}
vars-count 1  S ()  D (((_0 . closure)))  A ((_0 . μkanren-tag/sym98))  T ()
vars-count 1  S ()  D (((_0 . closure)))  A ((_0 . μkanren-tag/num99))  T ()
((λ (α) (assert (every (μ v (number? v)) (list α))) α))
\end{answers}
For a symbol the prefix is meaningful ($q \ne \cst{closure}$); for a number it is vacuous, and the
answer does not show it. A new store that reads D directly must expect such leftovers.

\subsection{The post-unification chain}

The goal \scm{=°} (\texttt{:539–552}, re-indented to fit the page) is where everything meets:
\begin{scheme}
(define ((=° u v) s)
  (μkanren-state-match ((vc S D A T tags) s)
    (define post°
      (λ1-match/first
        ((,s* ,assocs)
         (cond
           ((μkanren-state-equal? s s*) (✓° s*))          ; nothing bound: done
           ; D, touched prefixes only
           ((μkanren-verify-D/touched D s* assocs)
            => (μ D* (cond                                 ; then A, T, subsumptions
                       ((μkanren-verify-D/post D* A T s*) => ✓°)
                       (else (✗° s*)))))
           (else (✗° s*))))))
    (cond
      ((μkanren-state-unify u v s '()) => post°)
      (else (✗° s)))))
\end{scheme}
After a unification that binds something, the stores are re-verified in a hard-wired order:
\begin{center}
D (touched prefixes only) $\to$ A (in full) $\to$ D subsumed by A $\to$ T (in full) $\to$ D subsumed
by T $\to$ absentos of typed variables turned into D,
\end{center}
\begingroup\emergencystretch=1.5em
through \scm{μkanren-verify-D/post} (\texttt{:450–453}), \scm{μkanren-verify-A/post}
(\texttt{:424–\allowbreak 429}), \scm{μkanren-verify-T/post} (\texttt{:373}) and \scm{μkanren-subsume-T}
(\texttt{:352–356}), which rebuilds the state with the new D, A and T. Only D is incremental; A and T
are re-verified in full after every binding \scm{=°}. The whole chain is skipped only when nothing was
bound (\cref{prop:micro-eq}).\par\endgroup

\begin{kanren}[A new store has nowhere to plug in]
There is no list of verifiers to which a new store could append itself: the chain is a fixed sequence
of unexported procedures. λProlog features that need state of their own all hit this wall: a dynamic
program for AUGMENT \book{39}, scope information for GENERIC, a store of postponed non-pattern
equations that must be woken after each binding (\cref{ch:hounif}). What such a store requires of the
engine is (i) a place in the state that every goal copies forward, (ii) a call after each binding with
the new \scm{assocs}, (iii) a say in reification. \Cref{sec:micro-noedit} shows how far one gets
without (ii) and (iii); \cref{ch:design} proposes how to provide them.
\end{kanren}

% ------------------------------------------------------------------------------------------------
\section{Goals and the interleaving search}\label{sec:micro-search}

\subsection{Streams and goals}

\begin{definition}[Streams and goals]\label{def:micro-goal}
A stream (\texttt{(aux stream)}) is \scm{'()}, a \emph{mature} pair \scm{(state . stream)}, or an
\emph{immature} CHICKEN promise made with \scm{δ} (\scm{delay}, memoising). A \emph{goal} is a procedure
from a state to a stream of states. A \emph{relation} is a procedure from terms to a goal.
\end{definition}

The primitives (\texttt{:529–570}; an excerpt, comments ours) are few:
\begin{scheme}
; succeed once with the given state
(define ✓° list)
(define ✗° (K '()))                                    ; fail
(define ((freshª f) s)
  (μkanren-state-match ((vc S D A T tags) s)
    (let* ((α   (make-μkanren-var vc))
           (g   (f α))                                 ; the body is built only now
           (s*  (make-μkanren-state (add1 vc) S D A T tags)))
      (δ (g s*)))))                                    ; and suspended
(define ((orª f g) s) (append§/interleaved/2 (δ (f s)) (δ (g s))))
(define ((andª f g) s) (append-map§ g (δ (f s))))
\end{scheme}
The suffix \scm{ª} marks the \emph{applicative} procedure behind a macro; \scm{°} marks a goal, a
goal-making macro or a relation. Interleaving lives in \texttt{(aux stream)} (\texttt{aux.stream.scm:80–85}):
\begin{scheme}
(define append§/interleaved/2
  (λ (s1 s2)
    (cond
      ((null? s1) s2)
      ; SWAP on suspension
      ((promise? s1) (δ (append§/interleaved/2 s2 (force s1))))
      (else (cons§ (car s1) (append§/interleaved/2 (cdr s1) s2))))))
\end{scheme}
When the left stream is suspended, the two streams swap places; this is what makes disjunction
complete, and every \scm{δ} is a point where control can pass to another branch.

\paragraph{The macros} (\texttt{:575–593}, \texttt{:626–629}). \scm{(fresh° (x y) g ...)} nests one
\scm{freshª} per variable, each suspending once; \scm{(fresh° q ← (a b) g ...)} also binds \scm{q} to
\scm{(list a b)}. \scm{and°} is \emph{left}-nested,
\scm{(and° g1 g2 g3)} $=$ \scm{(andª (andª g1 g2) g3)}; \scm{or°} is \emph{right}-nested,
\scm{(or° g1 g2 g3)} $=$ \scm{(orª g1 (orª g2 g3))}. \scm{cond°} is \texttt{conde}:
\scm{(or° (and° g ...) ...)}, with no \texttt{else}. \scm{null°}, \scm{boolean°}, \scm{cons°} are
one-liners over \scm{=°} and \scm{or°}.

\begin{tutorial}[Fairness, and its shape]
Two infinite generators in a disjunction share the answers:
\begin{scheme}
(define-relation (nat° n) (or° (=° n 'z) (fresh° (m) (=° n `(s ,m)) (nat° m))))
(write (°->list/ground (take° 6 (fresh° (q) (or° (fresh° (n) (nat° n) (=° q `(left ,n)))
                                                (fresh° (n) (nat° n) (=° q `(right ,n)))))))) (newline)
\end{scheme}
\begin{answers}
((left z) (right z) (left (s z)) (right (s z)) (left (s (s z))) (right (s (s z))))
\end{answers}
The interleaving is fair between two branches, but \scm{or°} nests to the right, so with three
branches the first one competes against the \emph{pair} of the other two:
\begin{scheme}
(define-relation (numb° n k) (or° (=° n k) (numb° n (+ k 1))))
(write (°->list/ground (take° 12 (fresh° (q) (fresh° (n)
          (or° (and° (numb° n 0) (=° q (list 'a n)))
               (and° (numb° n 0) (=° q (list 'b n)))
               (and° (numb° n 0) (=° q (list 'c n))))))))) (newline)
\end{scheme}
\begin{answers}
((a 0) (a 1) (b 0) (a 2) (c 0) (a 3) (b 1) (a 4) (c 1) (a 5) (b 2) (a 6))
\end{answers}
In the long run branch \scm{a} gets half of the answers, \scm{b} and \scm{c} a quarter each (here 7, 3
and 2 of the first 12). Conjunction, on the other hand, is not fair:
\scm{(and° generator failing-test)} diverges, so as in every miniKanren generators go before tests and recursive calls last.
For instance, the \scm{append°} of \cref{sec:micro-reify} with its recursive call placed \emph{before}
\scm{(cons° a c rs)} still finds the four splits of \scm{(1 2 3)}, but then searches forever for a fifth
one under an unbounded run, while the order shown there terminates.
\end{tutorial}

\subsection{\texttt{define-relation} and the inverse-η delay}

Because \scm{and°} and \scm{or°} expand into \emph{procedure calls} of \scm{andª} and \scm{orª}, their
goal arguments are evaluated when the goal tree is built, before any state exists. A recursive relation
written with plain \scm{define} whose recursive call is not under a \scm{fresh°} therefore builds an
infinite tree:
\begin{scheme}
(define (ones° x) (or° (=° x 1) (ones° x)))          ; plain define: no delay
(write (°->list/ground (take° 3 (fresh° (q) (ones° q))))) (newline)
\end{scheme}
Run under \texttt{timeout 10}, this prints nothing and is killed (\texttt{exit=124}). The macro of
\texttt{:690} applies an inverse-η expansion and a suspension:
\begin{scheme}
(define-syntax-rule (define-relation (name arg ...) g ...) (define (name arg ...) (λ (s) (δ ((and° g ...) s)))))
\end{scheme}
so the body is built only when the goal meets a state, and then suspended:
\begin{scheme}
(define-relation (ones° x) (or° (=° x 1) (ones° x)))
(write (°->list/ground (take° 3 (fresh° (q) (ones° q))))) (newline)
\end{scheme}
\begin{answers}
(1 1 1)
\end{answers}
Plain \scm{define} remains fine when every recursive call sits under a \scm{fresh°}, whose body is under
\scm{λ (α)}; the relational interpreter of \texttt{test/microkanren-untagged.scm} is written that way.

\begin{pitfall}[Answer order is observable]
Each \scm{fresh°} variable and each relation call adds a suspension, \scm{or°} adds two. The test
suites pin the raw answers, \emph{including their order}. A new goal that adds or removes a single
\scm{δ} (for example a new \scm{=°}-like goal) can permute the answers of existing suites.
\end{pitfall}

\subsection{Soft cut, bounded runs, escape to Scheme}

\leavevmode\scm{(if° g? gt gf)} (\texttt{:596–603}) is LogicT's \texttt{ifte}, i.e.\ miniKanren's \texttt{conda}
for one clause: if \scm{g?} has at least one answer, \emph{all} its answers are continued with
\scm{gt}; otherwise \scm{gf} runs on the original state. \scm{(take° n g)} (\texttt{:605}) keeps the first
$n$ answers of \scm{g} (all of them if $n<0$); \scm{(take° 1 g)} is \texttt{once}. With \scm{gt} $=$ \scm{✗°}
and \scm{gf} $=$ \scm{✓°}, \scm{if°} is negation as failure:
\begin{scheme}
(write (°->list/ground (fresh° (q) (if° (or° (=° q 1) (=° q 2)) ✓° (=° q 3)))))
(newline)
(write (°->list/ground (fresh° (q) (if° ✗° ✓° (=° q 3))))) (newline)
(write (list (°->list/ground (fresh° (q) (=° q 1) (if° (=° q 2) ✗° ✓°)))
             (°->list/ground (fresh° (q) (=° q 2) (if° (=° q 2) ✗° ✓°)))
             (°->list/ground (fresh° (q) (if° (=° q 2) ✗° ✓°))))) (newline)
(write (°->list/ground (fresh° (q) (take° 2 (nat° q))))) (newline)
\end{scheme}
\begin{answers}
(1 2)
(3)
((1) () ())
(z (s z))
\end{answers}
The third line exhibits the classic unsoundness: with \scm{q} unbound, ``\scm{q} is not 2'' fails,
because the test succeeds by binding \scm{q}.

λProlog's \lp{if} and \lp{not} are defined with the cut \book{137}:
\begin{lprolog}
type if   o -> o -> o -> o.
if P Q R :- P, !, Q.
if P Q R :- R.

type not  o -> o.
not P :- P, !, fail.
not P.
\end{lprolog}
and the book's example of their non-logical character, ``\lp{?- X = 2, not (1 = X).} succeeds, whereas
\lp{?- not (1 = X), X = 2.} fails'' \book{137--138}, behaves the same way with \scm{if°}. There is one
difference: the cut keeps only the \emph{first} proof of \lp{P}, while \scm{if°} keeps all the answers
of its test, so λProlog's \lp{if P Q R} is \scm{(if° (take° 1 P) Q R)}:
\begin{scheme}
; X = 2, not (1 = X)
(write (list (°->list/ground (fresh° (q) (=° q 2) (if° (=° 1 q) ✗° ✓°)))
             ; the converse
             (°->list/ground (fresh° (q) (if° (=° 1 q) ✗° ✓°) (=° q 2)))))
(newline)
(write (list (°->list/ground (fresh° (q) (if° (or° (=° q 1) (=° q 2)) ✓° ✗°)))
             (°->list/ground
               (fresh° (q) (if° (take° 1 (or° (=° q 1) (=° q 2))) ✓° ✗°)))))
(newline)
\end{scheme}
\begin{answers}
((2) ())
((1 2) (1))
\end{answers}

\scm{(project° ((x* x) ...) g ...)} (\texttt{:626–627}) binds Scheme variables to the deep-walked values
of logic terms and runs \scm{g ...}: the non-relational escape hatch. An unbound variable is seen as a
\scm{μkanren-var} record.
\begin{scheme}
(write (°->list/ground (fresh° (q x) (=° x 20) (project° ((x* x)) (=° q (* 2 x*)))))) (newline)
(pretty-print (°->list #f (fresh° (q x) (project° ((x* x)) (=° q (list (μkanren-var? x*) (μkanren-var-index x*)))))))
\end{scheme}
\begin{answers}
(40)
((λ () (cons #t (cons 1 '()))))
\end{answers}

\begin{kanren}[AUGMENT and GENERIC as goals]
A goal sees the whole state and returns a stream of states, so it can do more than \scm{=°} does: it
can change the state on entry and post-process every answer state on exit. That is the shape that
both $D\imp G$ (add $D$, solve $G$, then restore the program on every answer of $G$) and $\forall x\,G$
(create $c$, solve $G$, then check that $c$ did not escape) require \book{39}. The two
prototypes of \cref{sec:micro-noedit} are written exactly this way, with \scm{append-map§} over the
answers of $G$.
\end{kanren}

% ------------------------------------------------------------------------------------------------
\section{Reification and the runners}\label{sec:micro-reify}

\subsection{From a state to a λ-expression}

An answer is not a substitution but a \emph{closed Scheme expression}
\begin{center}
\scm{(λ (α β ...) constraint-form ... term)}
\end{center}
which, applied to values for the free variables, re-checks the constraints with \scm{assert} and
\scm{deny} and returns the term. It is built by \scm{μkanren-project} (\texttt{:310–326}) in six steps:
\begin{enumerate}
\item deep-walk the query term;
\item name its free variables in order of first occurrence (left to right, car before cdr) by
  \emph{binding} each of them to a fresh reified variable $-1,-2,\dots$ in a scratch state
  (\scm{μkanren-state-reify}, \texttt{:207–222});
\item re-verify D in full;
\item \emph{purify}: drop every constraint that mentions a variable unreachable from the answer (still
  working in the scratch state); this is sound for disequalities, since such a variable can always be
  chosen to satisfy them;
\item drop the prefixes subsumed by others or made vacuous by a tag;
\item render with \scm{μkanren-form} (\texttt{:243–249}), sorting everything lexicographically.
\end{enumerate}

\begin{longtable}{@{}>{\raggedright\arraybackslash}p{0.3\fullwidthlen}>{\raggedright\arraybackslash}p{0.62\fullwidthlen}@{}}
\toprule
store & form in the answer \\
\midrule\endhead
D, prefix of one pair & \scm{(begin (deny (equal? α <repr>)))} \\*
D, prefix of $n$ pairs & \scm{(begin (deny (and (equal? α ...) (equal? β ...))))} \\*
A, grouped by tag & \scm{(assert (every (μ v (symbol? v)) (list α γ)))} \\*
T & \scm{(assert (absent? '<datum> <repr>))} \\
\bottomrule
\end{longtable}

\begin{scheme}
(pretty-print (°->list #f (fresh° (q x y) (=° q (list x y)) (≠° x 'a) (symbol° x) (absent° 'z y))))
(pretty-print (°->list #f (fresh° (q x y) (=° q (list x y)) (≠° (list x y) '(1 2)))))
(pretty-print (°->list #f (fresh° (q) (fresh° (h) (≠° h 3) (=° q 7)))))
(pretty-print (°->list/ground (fresh° (q x) (=° q (list x 1)) (≠° x 2))))
\end{scheme}
\begin{answers}
((λ (α β)
    (begin (deny (equal? α 'a)))
    (assert (every (μ v (symbol? v)) (list α)))
    (assert (absent? 'z β))
    (cons α (cons β '()))))
((λ (α β)
    (begin (deny (and (equal? α 1) (equal? β 2))))
    (cons α (cons β '()))))
((λ () 7))
((α 1))
\end{answers}
The third answer shows purification: the disequality on the unreachable \scm{h} is gone. The fourth
shows grounding: \scm{°->list/ground} evaluates only \scm{(lambda vars term)}, \emph{dropping the
constraints}, and applies it to the variable symbols themselves, so the free variable reads back as the
symbol \scm{α}. The show module \texttt{(aux kanren micro show)} parses the answer forms back without
\scm{eval} (\texttt{aux.kanren.micro.show.scm:58–72}), so a new kind of constraint in answers needs a
case there too.

\subsection{Running}

The runners are at \texttt{:686–713}:
\begin{scheme}
(define (°->§ g)
  (let* ((§ (δ (g μkanren-state-empty)))
         (P (μkanren-project (make-μkanren-var 0))))
    (map§ P §)))
(define-syntax-rule (μkanren-run (v n grounded) g ...) (°->list grounded (take° n (fresh° (v) g ...))))
\end{scheme}
\scm{(°->list grounded g)} forces \scm{°->§} and, if \scm{grounded}, \scm{eval}s each answer in micro's
own module environment; \scm{°->list/ground} is \scm{(°->list #t g)}. \scm{μkanren-run} is the
TRS-like \texttt{run}, with $n<0$ meaning ``all''. The query variable is always \emph{the variable of
index 0}, i.e.\ the first one the goal introduces. With the usual \scm{append°}:
\begin{scheme}
(define-relation (append° r s rs)
  (cond°
    ((null° r) (=° s rs))
    ((fresh° (a d c) (cons° a d r) (cons° a c rs) (append° d s c)))))
(pretty-print (μkanren-run (q -1 #t) (fresh° (x y) (append° x y '(1 2 3)) (=° q (list x y)))))
(pretty-print (μkanren-run (q 3 #f) (fresh° (x y z) (append° x y z) (=° q (list x y z)))))
\end{scheme}
\begin{answers}
((() (1 2 3)) ((1) (2 3)) ((1 2) (3)) ((1 2 3) ()))
((λ (α) (cons '() (cons α (cons α '()))))
 (λ (α β) (cons (cons α '()) (cons β (cons (cons α β) '()))))
 (λ (α β γ)
    (cons (cons α (cons β '())) (cons γ (cons (cons α (cons β γ)) '())))))
\end{answers}

\begin{pitfall}[The query variable is index 0]
\scm{°->§} hard-wires the empty initial state and projects variable 0. An extension that reserves
indices (for eigenvariables, registers, \dots) \emph{before} the first \scm{fresh°} shifts the query
variable; a module that needs a richer initial state needs its own runner, which is easy since
\scm{make-μkanren-var} and \scm{μkanren-project} are exported.
\end{pitfall}

\begin{kanren}[What answers will have to show]
Eigenvariables and a dynamic program never appear in today's answers. If a \scm{≠°} mentioning an
eigenvariable survives its scope, or a postponed higher-order equation is left at the end, reification
must either drop it or print it; purification (step 4) does not know about scopes. See
\cref{sec:micro-forall} for an example and \cref{ch:design} for the proposal.
\end{kanren}

% ------------------------------------------------------------------------------------------------
\section{Aggregation}\label{sec:micro-aggr}

Four SQL-like goals (\texttt{:631–684}) consume the \emph{whole} answer stream of a sub-goal with
\scm{foldr§} and turn it into Scheme data, so the sub-goal must be finite and the goals are
non-relational (\texttt{findall}/\texttt{bagof}-like). \scm{groupby°} groups the answers of \scm{g} by
the values of the keys and runs the continuation once per group \emph{on the original state}, with the
aggregate and the key names bound as Scheme variables; \scm{window°} is \texttt{OVER (PARTITION BY ...)}
and runs the continuation on each answer state of \scm{g}; \scm{set°} folds per key; \scm{enumerate°}
numbers the answers in stream order.
\begin{scheme}
(pretty-print (°->list/ground (fresh° (q) (fresh° (k v)
   (groupby° (((vs (λ (l) l)) v)) over (k) from
             (or° (and° (=° k 'a) (=° v 1)) (and° (=° k 'b) (=° v 2)) (and° (=° k 'a) (=° v 3)))
             => (=° q (list k vs)))))))
\end{scheme}
\begin{answers}
((b (2)) (a (1 3)))
\end{answers}
They are the precedent for goals that run a sub-search to completion, the mechanism a
\texttt{findall}-style predicate would use.

% ------------------------------------------------------------------------------------------------
\section{Extension points}\label{sec:micro-ext}

\subsection{Open versus closed}

\Cref{tab:micro-open} sorts the engine's parts by whether an importing module can use or extend them
without editing \texttt{aux.kanren.micro.scm}.

\begin{longtable}{@{}>{\raggedright\arraybackslash}p{0.47\fullwidthlen}>{\raggedright\arraybackslash}p{0.47\fullwidthlen}@{}}
\caption{The inventory of \micro{} for an extension.}\label{tab:micro-open}\\
\toprule
\textbf{open} (usable from outside) & \textbf{closed} (needs an edit of micro) \\
\midrule\endfirsthead
\toprule
\textbf{open} & \textbf{closed} \\
\midrule\endhead
Goals are plain procedures \emph{state} $\to$ \emph{stream}; all field accessors and
\scm{μkanren-state-with} are exported.
& The field list of the state (\texttt{:69}), its positional constructions (\texttt{:71}, \texttt{:83},
\texttt{:132}, \texttt{:536}), the six-ary \scm{μkanren-state-match} and its six users. \\
Raw unification \scm{μkanren-state-unify} returns the new state and the exact new bindings (but does
not re-verify D/A/T).
& The unification \scm{cond} (\texttt{:160–181}): no hook for a term kind with non-syntactic equality. \\
Walks \scm{μkanren-state-find}, \scm{/value}, \scm{/values}, \scm{/repr}.
& The binding site with occurs check, \scm{μkanren-state-update}, unexported: the only place to refuse a binding before it enters the state (a check right after each
\scm{=°} works as well, \cref{sec:design-hooks}). \\
Variables: \scm{make-μkanren-var}, \scm{μkanren-var-index}, \dots; a fresh unique integer by bumping
\texttt{vars-count}.
& The post-\scm{=°} chain (\texttt{:539–552}, \texttt{:373–478}), wired for D/A/T, its pieces unexported. \\
New first-order constructors: any \scm{define-record} unifies, reifies, works with \scm{≠°} and
\scm{absent°}.
& Reification (\texttt{:310–326}, \texttt{:243–249}): knows only D/A/T; purification and subsumption
are store-specific. \\
New disjoint types: \scm{define-μkanren-tag} with \scm{μkanren-make-tag-A}.
& \scm{°->§}: fixed empty initial state, fixed query index 0. \\
Projection \scm{μkanren-project}; the combinators \scm{freshª}, \scm{andª}, \scm{orª}, \scm{if°},
\scm{take°}, \scm{project°}; all of \texttt{(aux stream)}; \scm{define-relation}. & \\
\bottomrule
\end{longtable}

Two of the open entries in action: a user-defined disjoint type, and a goal written as a bare lambda
that reads the state.
\begin{scheme}
(define-μkanren-tag (tag/str v) (string? v))
(define string° (μkanren-make-tag-A tag/str))
(pretty-print (°->list #f (fresh° (q) (string° q))))
(pretty-print (map (λ (g) (°->list #f g))
                   (list (fresh° (q) (string° q) (=° q "hi"))
                         (fresh° (q) (string° q) (=° q 3))
                         (fresh° (q) (string° q) (symbol° q))
                         (fresh° (q) (string° q) (≠° q 3)))))
(pretty-print (°->list/ground (fresh° (q x y) (λ (s) ((=° q (μkanren-state-vars-count s)) s)))))
\end{scheme}
\begin{answers}
((λ (α) (assert (every (μ v (string? v)) (list α))) α))
(((λ () "hi")) () () ((λ (α) (assert (every (μ v (string? v)) (list α))) α)))
(3)
\end{answers}
The new tag is printed in answers, is disjoint from \scm{symbol°}, and subsumes \scm{(≠° q 3)}. The
last goal sees \texttt{vars-count} $=3$.

A CHICKEN rule constrains every new module: a macro exported from a module whose expansion refers to
an unexported binding fails at the use site. This is why \scm{freshª}, \scm{andª}, \scm{orª},
\scm{✓°}, \scm{✗°}, \scm{=°} and \scm{μkanren-group§} are exported (\texttt{:4–17}); macros such as
a future \texttt{∀°} or \texttt{⊃°} must follow the same discipline.

\subsection{Three tricks that need no edit, and their limits}\label{sec:micro-noedit}

\paragraph{Trick 1: wrapper goals that emit \texttt{=°}.} A solver written outside the engine (say,
higher-order pattern unification computing a most general unifier $\{F\mapsto t_1,\dots\}$) should
express its result as ordinary \scm{=°} goals, \scm{(and° (=° F t1) ...)}, so that each binding runs
the D/A/T chain. The tempting shortcut, calling the exported raw unification and returning its state,
silently bypasses the stores:
\begin{scheme}
; WRONG: bypasses the D/A/T re-verification
(define ((bind/raw x t) s)
  (let1 (r (μkanren-state-unify x t s '())) (if r (✓° (car r)) (✗° s))))
; right: the result is expressed as an ordinary =°
(define ((bind/emit x t) s) ((=° x t) s))

(define raw (§->list ((fresh° (q) (≠° q 'a) (bind/raw q 'a)) μkanren-state-empty)))
; q = a and q ≠ a!
(pretty-print (map (λ (s) (list (sbral->list (μkanren-state-S s)) (μkanren-state-D s))) raw))
(pretty-print (handle-exceptions e (list 'error ((condition-property-accessor 'exn 'message) e))
                (°->list #f (fresh° (q) (≠° q 'a) (bind/raw q 'a)))))
(pretty-print (°->list #f (fresh° (q) (≠° q 'a) (bind/emit q 'a))))
\end{scheme}
\begin{answers}
(((a) (((_0 . a)))))
(error "bad argument type - not a list")
()
\end{answers}
The raw version produces a state where S says $q=\cst{a}$ and D says $q\ne\cst{a}$; reification then
re-verifies D, gets \scm{#f} for ``violated'', and crashes. The emitting version fails, as it should.
\emph{Limits}: extra suspensions and repeated verification; and plain \scm{=°} applied by a user to
terms that need the special solver would still unify them syntactically, so a module must forbid that
or shadow \scm{=°}.

\paragraph{Trick 2: S-registers.} There is no way to add a field from outside: every goal of micro
rebuilds the state positionally, so a wrapper record would be lost at the next \scm{=°}. Dynamic
scoping (\scm{parameterize}) fails too, because promises are forced outside its extent, and so does a
side table keyed by the state, because a new state is allocated at every binding. But S itself is
copied forward by every goal and is persistent and branch-local. So: reserve a fresh index $n$ (bump
\texttt{vars-count}), bind variable $n$ by raw unification to an opaque record carrying the data, and
find the current value as the register of highest index. Raw unification is harmless here because the
register variable is fresh and unconstrained, and registers are never reachable from the query variable,
so they are never reified. Here the data is a λProlog-style program context, giving a prototype of
AUGMENT:
\begin{scheme}
; a register: an opaque record bound in S
(define-record ctx clauses)
; reserve index n, bind var n (raw unify) to a ctx record
(define (ctx-push s Γ)
  (let* ((n (μkanren-state-vars-count s))
         (s1 (μkanren-state-with s vars-count: (add1 n))))
    (car (μkanren-state-unify (make-μkanren-var n) (make-ctx Γ) s1 '()))))
; the most recent register, by a shallow walk downwards
(define (ctx-current s)
  (let loop ((i (sub1 (μkanren-state-vars-count s))))
    (cond ((< i 0) '())
          (else (let1 (v (μkanren-state-find (make-μkanren-var i) s))
                  (if (ctx? v) (ctx-clauses v) (loop (sub1 i))))))))
; D ⊃ G, D a list of clause procedures (λ (x) goal)
(define ((⊃ª D g) s)
  (let1 (Γ (ctx-current s))
    (append-map§ (λ (s2) (✓° (ctx-push s2 Γ))) (δ (g (ctx-push s (append D Γ)))))))
; a dynamic predicate: try every clause of the context
(define ((hyp° x) s)
  (let1 (Γ (ctx-current s)) ((fold-right (λ (c g) (orª (c x) g)) ✗° Γ) s)))

(define (is° v) (λ (x) (=° x v)))            ; the clause  hyp X :- X = v
(for-each (λ (g) (write (°->list/ground g)) (newline))
  (list (fresh° (q) (⊃ª (list (is° 'a)) (hyp° q)))
        (fresh° (q) (⊃ª (list (is° 'a)) ✓°) (hyp° q))
        (fresh° (q) (⊃ª (list (is° 'a)) (⊃ª (list (is° 'b)) (hyp° q))))
        (fresh° (q) (or° (⊃ª (list (is° 'a)) (hyp° q)) (hyp° q)))
        (fresh° (q) (fresh° (y) (⊃ª (list (λ (x) (=° x y))) (fresh° (u) (hyp° u) (=° u 7))) (=° q y)))))
\end{scheme}
\begin{answers}
(a)
()
(b a)
(a)
(7)
\end{answers}
In order: the assumed clause is used; outside its scope it is gone; nested assumptions stack, newest
first; two branches of a disjunction do not see each other's assumptions; and a free logic variable
\scm{y} of an assumed clause is \emph{shared} with the context, not renamed, as in λProlog
\book{86--88}. \emph{Limits}: the lookup scans the indices, $O(\texttt{vars-count}\cdot\log)$ per call;
registers must be allocated after the query variable; it is a library convention, not an engine
guarantee.

\subsection{Trick 3: a \texttt{∀} prototype with a post-filter scope check}\label{sec:micro-forall}

The eigenvariable is a record whose id is a reserved index, hence unique; distinct eigenvariables never
unify (their id slots differ). The scope check is a post-filter on each answer state of the body: every
variable of index $< n$, i.e.\ older than the eigenvariable, is deep-walked and must not contain it.
\begin{scheme}
(define-record eigen id)
(define (contains? c t)
  (cond
    ; find/value REBUILDS records: compare ids, not identity
    ((eigen? t) (= (eigen-id t) (eigen-id c)))
    ((pair? t) (or (contains? c (car t)) (contains? c (cdr t))))
    ((vector? t) (any (λ (u) (contains? c u)) (vector->list t)))
    ((and (record-instance? t) (not (μkanren-var? t)))
     (any (λ (u) (contains? c u)) (vector->list (record->vector t))))
    (else #f)))
(define ((∀ª f) s)
  (let* ((n (μkanren-state-vars-count s))
         (c (make-eigen n))                        ; unique: index n is reserved
         (s* (μkanren-state-with s vars-count: (add1 n)))
         (escapes? (λ (s1) (let loop ((i 0))
                             (and (< i n)
                                  (or (contains? c (μkanren-state-find/value (make-μkanren-var i) s1))
                                      (loop (add1 i))))))))
    (append-map§ (λ (s1) (if (escapes? s1) '() (✓° s1))) (δ ((f c) s*)))))
(define-syntax-rule (∀° (x) g ...) (∀ª (λ (x) (and° g ...))))

(for-each (λ (g) (pretty-print (°->list #f g)))
  (list (fresh° (q) (∀° (x) (=° q x)))                              ; 1. ∃Q ∀x. Q = x
        (fresh° (q) (∀° (x) (fresh° (y) (=° y x))) (=° q 'ok))      ; 2. ∀x ∃Y. Y = x
        (fresh° (q) (∀° (x) (∀° (y) (=° x y))))                     ; 3. ∀x ∀y. x = y
        (fresh° (q) (∀° (x) (fresh° (y) (=° q (list y)) (=° y x)))) ; 4. alias
        (fresh° (q) (∀° (x) (fresh° (z) (=° q (list 'f z)))))       ; 5.
        (fresh° (q) (∀° (x) (≠° q x)))))                            ; 6. ≠ survives
\end{scheme}
\begin{answers}
()
((λ () 'ok))
()
()
((λ (α) (cons 'f (cons α '()))))
((λ (α) (begin (deny (equal? α (make-record-instance 'eigen 1)))) α))
\end{answers}

\begin{tutorial}[Reading the six answers]
\begin{enumerate}
\item $\exists Q\,\forall x.\,Q=x$ is the book's \lp{sigma x\ pi y\ x = y}, ``not provable'' \book{85}.
  Here \scm{q} is \scm{_0}; \scm{∀°} reserves index 1 and makes \scm{c} $=$ \scm{(eigen 1)}; the body
  binds \scm{_0 := c}; the post-filter walks \scm{_0}, finds \scm{c}, and discards the state.
\item $\forall x\,\exists Y.\,Y=x$ succeeds: \scm{y} has index 2 $\ge n=1$, so it is inside the scope.
\item Two eigenvariables \scm{(eigen 1)} and \scm{(eigen 2)} differ in their id slot: unification fails.
\item The escape is indirect: \scm{_0 := (list _2)} and \scm{_2 := c}. Because the check deep-walks the
  old variable, it sees \scm{(list c)} and rejects the state. This is the chain created by the
  older-to-newer orientation of \cref{sec:micro-unify}.
\item A younger variable left unbound is harmless and is reified as \scm{α}.
\item A disequality mentioning \scm{c} survives in D and is reified with the eigenvariable in it: the
  check filters bindings, not constraints.
\end{enumerate}
\end{tutorial}

Why is checking at the end of the body enough for soundness? After the scope closes, \scm{c} is
lexically unreachable and, by the check, no term reachable from an older variable contains it, so no
later goal can bring it back.\sidenote{This argument is informal; the prototype was tested, not proved.} What
is lost is \emph{early pruning}: a bad binding is detected only when the body has produced an answer,
at a cost of $O(n\cdot|\text{term}|)$ per answer. Eager checking needs the scope test at the binding
site (closed), or a post-unification hook with a level map, or a module-level \scm{=°} that shadows
micro's and checks its bindings against the open scopes kept in an S-register; the last option covers
only goals written with the new module's \scm{=°}.

\begin{pitfall}[The scope check must not use \texttt{eq?}]
With the first line of \scm{contains?} written as \scm{((eigen? t) (eq? t c))}, the same query as
answer~1 wrongly succeeds, because \scm{find/value} rebuilt the record (\cref{sec:micro-subst}):
\begin{scheme}
(pretty-print (°->list #f (fresh° (q) (∀° (x) (=° q x)))))
\end{scheme}
\begin{answers}
((λ () (make-record-instance 'eigen 1)))
\end{answers}
\end{pitfall}

\begin{kanren}[What the tricks leave open]
The three tricks show that the shape of AUGMENT and GENERIC fits \micro{} today: a goal that changes
the state on entry and filters or restores it on every answer. What they cannot provide is what only
the engine can do: a \emph{named} place in the state (instead of a register scan), a call after every
binding (for eager scope checks and for waking postponed equations), and a say in reification (to drop
out-of-scope constraints and print new ones). These are the requirements \cref{ch:design} turns into a
proposal, and \cref{ch:module} into a module.
\end{kanren}

% ------------------------------------------------------------------------------------------------
\section{Summary}\label{sec:micro-summary}

\begin{itemize}
\item Terms are Scheme data; logic variables are records with an integer index, working ($\ge 0$,
  printed \scm{_i}) or reified ($<0$, printed \scm{α}, \scm{β}, \dots). Indices grow monotonically along a
  branch, so an index is an age and doubles as a level boundary.
\item A state has six fields \texttt{vars-count}, S, D, A, T, \texttt{tags} (the last one dead). States
  are compared by \scm{eq?}: unification returns the same state iff it binds nothing, and every
  entailment test relies on this.
\item S is an sbral, most recent variable first, extended lazily; the walk follows triangular bindings
  and treats reified variables as constants; deep walks rebuild records, losing their identity.
\item Unification is first-order with occurs check; the only binding site is
  \scm{μkanren-state-update}; the older variable is bound to the newer; vectors and records unify as
  slot lists (so records and same-shaped vectors unify); numbers by a NaN-reflexive \scm{eqv?}; closures
  only by \scm{eq?}.
\item D holds prefixes (negated conjunctions of bindings), re-verified incrementally on the touched
  prefixes; A holds disjoint type tags, T absento constraints, both re-verified in full; the chain after
  \scm{=°} is fixed and has no slot for a new store.
\item Goals are \emph{state} $\to$ \emph{stream}; \scm{orª} interleaves by swapping on suspensions
  (\scm{or°} nests right, so three branches get $\frac12,\frac14,\frac14$); conjunction is not fair;
  \scm{define-relation} delays recursive relations; \scm{if°} is \texttt{conda} (it keeps every answer
  of its test, so λProlog's cut-based \lp{if} is \scm{if°} on a \scm{(take° 1 ...)} test), \scm{take°} bounds a run,
  \scm{project°} escapes to Scheme; answer order is pinned by the tests.
\item Answers are closed expressions \scm{(λ (vars) constraints ... term)}, purified of unreachable
  variables; grounding drops the constraints; the query variable is index 0.
\item Open to extension: goals as procedures, raw unification, walks, records as constructors, user
  type tags, projection, the combinators. Closed: the state's fields, the unification \scm{cond}, the
  binding site, the post-\scm{=°} chain, reification, the runner.
\item Without editing micro one can: emit \scm{=°} from external solvers (never return raw-unified
  states), carry branch-local data in S-registers (a working AUGMENT prototype), and check
  eigenvariable scope by post-filtering the answers of a \scm{∀} body (a working GENERIC prototype). Their
  limits — scans, late pruning, constraints outliving scopes — are the requirements of \cref{ch:design}.
\end{itemize}
